\section{Preliminaries}
% Introducción a la idea
The nature of quantum measurement is intrinsically probabilistic, with this in mind, we are interested in studying probabilistic evolution for lambda calculi. This evolution can be thought as a non-deterministic reduction step, with every option is weighted with a probability where every option sums to 1.

This inclusion adds a layer of complexity to the usual mechanisms over traditional rewrite systems. For example, properties like strong termination or confluence take on different definitions to accomodate this characteristic. We will focus particularly on confluence of probabilistic calculi for our first chapter.

% Definición de confluencia tradicional
To give a simple overview, confluence (otherwise known as the Church-Rosser property) is a property that states that no matter which redex is reduced on a term $t$ 

% Breve explicación de porque no funciona

When dealing with probabilistic lambda calculus, we can find two different sources of divergence.
\begin{itemize}
\item A \textit{single redex} may reduce in two different ways via a probabilistic reduction.

\item A term with \textit{multiple redexes} and no strategy, could be reduced in different ways.
\end{itemize}
For example, we can consider a lambda calculus extended with a coin $\coin$ reducing to $0$ or $1$ with probability $\frac 12$ each. Then, taking just the coin, we are in the first case of divergence. While taking, for example, $(\lambda x.\lambda y.yxx)\coin$, we are in the second case, since we can either beta reduce, or reduce the coin.

There is no point in trying to achieve confluence in the first case: the coin is non-confluent by design. However, we can analyse the branching paths and verify that the probability of reducing to a particular term stays the same, regardless of the reduction sequence. This is what we call \textit{probabilistic confluence}.

To study this kind of cases, Bournez and Kirchner developed the notion of PARS~\cite{BK02}, later refined in \cite{BG06}. Using the techniques described in~\cite{DCM17,Faggian} we can define the rewriting rules over the distributions. 

% Vamos a tener que trabajar con PARS (Definición)

If we denote by $[(p_1,t_1),\dots,(p_n,t_n)]$ the probability distribution where $t_i$ has probability $p_i$, the possible reductions from the previous example are depicted in Figure~\ref{fig:noconfl}. The resulting distributions are not only different, but also divergent, since in the left branch, the probability to arrive, for example, to $\lambda y.y01$ is $0$, while it is one of the possible results in the right branch.

\begin{figure}[t]
    \centering
    \begin{tikzcd}[column sep=small]
    & (\lambda x.\lambda y.yxx)\coin \ar[dl]\ar[dr] & \\
    \left[\begin{array}{c}
        ({\frac 12}, (\lambda x.\lambda y.yxx)0),\\[5pt]
        ({\frac 12}, (\lambda x.\lambda y.yxx)1)
        \end{array}\right]\ar[d] 
    & &  
    {[(1, \lambda y.y\coin\coin)]}\ar[d,"*",pos=1]\\
    \left[\begin{array}{c}
        ({\frac 12}, \lambda y.y00),\\[5pt]
        ({\frac 12}, \lambda y.y11)
    \end{array}\right]
    &&\hspace{-50pt}
    \left[\begin{array}{c}
        ({\frac 14}, \lambda y.y00), 
        ({\frac 14}, \lambda y.y01),\\[5pt]
        ({\frac 14}, \lambda y.y10), 
        ({\frac 14}, \lambda y.y11)
    \end{array}\right]
    \end{tikzcd}
    \caption{Counterexample of confluence}
    \label{fig:noconfl}
\end{figure}

